\section{Discussion}\label{sec:discussion}
\subsection{Defects exposed by the evaluation}
Running the real system on real talks surfaced problems that the offline test suite, which
uses fake providers, could not reveal. We report them because they are the most practically
useful outcome for practitioners building similar pipelines.

\textbf{Guardrail scope (fixed before the final runs).} The guardrail originally saw only
the first 12{,}000 transcript characters and not the author profile. In the pilot, all six
drafts of a talk were flagged: user-supplied sign-offs were reported as ``invented'', and
correct facts from late in the talk (supplied to the writers by RAG) would have been
penalised. The guardrail now receives the whole transcript and the profile.

\textbf{Few-shot template leakage (fixed).} The flyer and summary skill files contained
example outputs with a concrete date (``October~24, 2026'') and a venue, and the flyer rules
demanded a date. With no event metadata available, writers copied the example date into
drafts; the guardrail caught it, but only after a regeneration round-trip. A shared rule
forbidding copying of example details and requiring ``TBA'' for unknown logistics removed
the behaviour (Table~\ref{tab:format}).

\textbf{Sticky model fallback (fixed after the evaluation).} The Vertex provider remembered the
last working model for the whole process. In the pilot, a single empty response from
\texttt{gemini-3.1-flash-lite} switched every subsequent call of the process to the larger
\texttt{gemini-3.6-flash}, silently changing quality, latency and cost. The provider now always
starts with the configured model and skips a model only while its own circuit breaker is open
(three consecutive failures, 60\,s). In the repeated fault injection the next call returns to
\texttt{gemini-3.1-flash-lite} after an empty response and after a persistent 429
(Table~\ref{tab:faults}).

\textbf{Inert circuit breaker (fixed after the evaluation).} The circuit breaker recorded
failures but its \texttt{allow()} method never raised, so it never short-circuited a call. It
now rejects calls while open and lets one trial call through after the window (half-open).
The probe now shows 3 of 8 consecutive failing calls short-circuited (0 before), and the
per-model breakers mean that a rate-limited primary costs the back-off delay only on the first
three calls (Section~\ref{sec:reliability}).

\textbf{Silent offline fallback (fixed after the evaluation).} When every model failed, the
provider returned offline template text and the job looked like a normal result. Every step
that falls back (offline template, placeholder draft, missing translation) now marks the job
as degraded; such a job ends in \texttt{needs\_review} with an explanatory message. All 20
fault-injection trials in which every model failed are flagged.

\textbf{Unguarded response shape (fixed after the evaluation).} In
\LiveAnalysisShapeFailures{} live runs---every run of one talk (\emph{Elizabeth Gilbert})---the
analysis model returned a JSON array instead of an object, the analysis node raised
\texttt{'list' object has no attribute 'get'} and the job failed. The failure is deterministic
for that input, so retries did not help. Replies are now coerced to an object (an array of
objects is unwrapped and merged), and the analysis node re-asks once with an explicit
instruction if the expected keys are still missing. With the fix all seven runs of the talk
completed, each with a single analysis call, and the talk is included in the ablation and
quality results.

\textbf{Regional endpoints matter.} Embedding calls on Vertex's \texttt{global} endpoint took
about 12\,s each from our location versus 0.5--0.7\,s on regional endpoints; moving
embeddings to \texttt{asia-south1} removed RAG indexing from the critical path.

\textbf{Network resilience of the host.} Several runs coincided with machine-wide network
outages (DNS failures, SSL errors, stalled connections). Because token refresh in the
Google client is synchronous, a stalled refresh blocks the event loop and the per-call
timeout cannot fire; one call stalled for over 80 minutes before completing. The fallback
STT engine (Gemini API) also failed in these cases because its API key was invalid, and the
runner did not fall back a second level. Credentials of every configured engine are now
checked at start-up, a rejected engine is skipped, and speech-to-text tries every remaining
engine in turn (fixed after the evaluation). Following these stalls we added a 240\,s per-chunk timeout to
the Vertex STT provider; the end-to-end runs used the version without it, the later
cached-stage transcription runs the version with it.

\subsection{Threats to validity}
\emph{Sample size}: \DatasetTalks{} talks and \DatasetLong{} long talks are enough for paired
non-parametric tests and descriptive statistics, not for fine-grained claims; confidence
intervals are wide. \emph{Domain}: TED talks are well-recorded, single-speaker English
talks; noisy multi-speaker event audio would raise WER. \emph{Automatic judging}: quality
scores come from an LLM judge, not human raters; we use a different, larger judge model,
blind conditions, a human reference transcript and a second judge, but judge bias towards
fluent text remains. Key points are model-extracted (with verbatim quotes checked
mechanically). \emph{Retrieval labels} are produced by construction from model-written
questions and may favour paraphrase-friendly queries. \emph{Environment}: latency was
measured on a consumer laptop and home network with two concurrent jobs; absolute numbers
depend on region and load. \emph{Cost}: we report measured tokens; we could not retrieve
official per-token prices for the models used, so monetary cost is not reported.
